\section{Étude de cas : Banque Atlas}

\subsection{6\_init — Initialisation du schéma BANQUE}
\textbf{Objectif.} Créer le schéma BANQUE (tables clients, comptes, opérations, vue v\_clients\_guichet,
procédure p\_operation) et peupler les données de test.

\shots{6_init}

\textbf{Observations et réponses.}
\begin{enumerate}
  \item L'utilisateur \code{BANQUE} est créé avec \code{NO AUTHENTICATION} et
        \code{QUOTA UNLIMITED ON users} (schéma seulement, pas de connexion directe).
  \item Les tables \code{clients} (4 lignes), \code{comptes} (5 lignes) et
        \code{operations} sont créées et peuplées.
  \item La vue \code{v\_clients\_guichet} expose \code{client\_id, nom, cin, telephone}
        (pas de \code{revenu\_mensuel}, pas de \code{conseiller}).
  \item La procédure \code{p\_operation} se compile sans erreur : elle est la seule porte
        d'entrée pour les mouvements d'argent (droits du définisseur), contrôle le solde et
        journalise dans \code{operations} avec \code{SYS\_CONTEXT('USERENV','SESSION\_USER')}.
\end{enumerate}
\ecart{Aucun.}

\subsection{Exercice 6.1 — T1 : Matrice RBAC}
\textbf{Objectif.} Construire la matrice rôles × objets et justifier chaque case au regard des
exigences E1 à E3.

\begin{center}
\small
\begin{tabular}{|l|c|c|c|c|c|c|}
\hline
\textbf{Rôle} & \code{CREATE SESSION} & \code{clients} & \code{comptes} & \code{operations} & \code{v\_clients\_guichet} & \code{p\_operation} \\
\hline
\code{R\_BQ\_CONNEXION} & \textbf{Oui} & -- & -- & -- & -- & -- \\
\hline
\code{R\_BQ\_GUICHET} & via CONNEXION & -- & SELECT & -- & SELECT & EXECUTE \\
\hline
\code{R\_BQ\_CONSEILLER} & via GUICHET & SELECT, INSERT & INSERT & SELECT & SELECT (hérité) & EXECUTE (hérité) \\
\hline
\code{R\_BQ\_DIRECTEUR} & via CONSEILLER & (hérité) & (hérité) & (hérité) & (hérité) & EXECUTE (hérité) \\
\hline
\code{R\_BQ\_AUDIT} & \textbf{Oui} & SELECT & SELECT & SELECT & -- & -- \\
\hline
\code{R\_BQ\_BATCH} & via CONNEXION & -- & SELECT & -- & -- & EXECUTE \\
\hline
\end{tabular}
\end{center}

\textbf{Héritages :} \code{R\_BQ\_GUICHET} $\to$ \code{R\_BQ\_CONNEXION} ;
\code{R\_BQ\_CONSEILLER} $\to$ \code{R\_BQ\_GUICHET} $\to$ \code{R\_BQ\_CONNEXION} ;
\code{R\_BQ\_DIRECTEUR} $\to$ \code{R\_BQ\_CONSEILLER} ; \code{R\_BQ\_AUDIT} et
\code{R\_BQ\_BATCH} $\to$ \code{R\_BQ\_CONNEXION}.

\textbf{Justification (E1–E3) :}
\begin{enumerate}
  \item \textbf{E1} (aucun privilège direct aux utilisateurs métier) : tous les privilèges
        sont accordés aux rôles, jamais directement aux comptes GUI1, GUI2, CONS1, DIR1,
        AUD1, APP\_BATCH.
  \item \textbf{E2} (pas de ANY, pas de DBA/CONNECT/RESOURCE direct) : aucun compte métier
        ne possède ces rôles/privilèges directement.
  \item \textbf{E3} (soldes modifiables uniquement par P\_OPERATION) : aucun rôle n'a
        \code{UPDATE} sur \code{comptes.solde} ; seule la procédure \code{p\_operation}
        (droits du définisseur) modifie les soldes.
\end{enumerate}
\textbf{Seul privilège avec retransmission :} \code{GRANT r\_bq\_guichet TO dir1 WITH ADMIN
OPTION} (accordé directement au compte DIR1, pas via le rôle R\_BQ\_DIRECTEUR).

\subsection{Exercice 6.2 — T2 : Profils}
\textbf{Objectif.} Créer \code{P\_BQ\_AGENT}, \code{P\_BQ\_CONSEIL}, \code{P\_BQ\_SENSIBLE}
et \code{P\_BQ\_SERVICE} conformément à E4 et E5.

\shots{6_2}

\textbf{Observations et réponses.}
\begin{enumerate}
  \item \code{P\_BQ\_AGENT} : \code{SESSIONS\_PER\_USER 1}, \code{IDLE\_TIME 10},
        \code{FAILED\_LOGIN\_ATTEMPTS 3}, \code{PASSWORD\_LOCK\_TIME 30/1440},
        \code{PASSWORD\_LIFE\_TIME 60}, \code{PASSWORD\_REUSE\_TIME 365},
        \code{PASSWORD\_REUSE\_MAX 5}, \code{PASSWORD\_VERIFY\_FUNCTION ORA12C\_VERIFY\_FUNCTION}.
  \item \code{P\_BQ\_CONSEIL} : identique à \code{P\_BQ\_AGENT} sauf
        \code{SESSIONS\_PER\_USER 2} (le conseiller a besoin de 2 sessions).
  \item \code{P\_BQ\_SENSIBLE} : mêmes règles E4 que l'agent, avec
        \code{ORA12C\_STRONG\_VERIFY\_FUNCTION} (E5 : mot de passe renforcé pour directeur
        et auditeur).
  \item \code{P\_BQ\_SERVICE} : \code{PASSWORD\_LIFE\_TIME UNLIMITED} (pas d'expiration),
        \code{IDLE\_TIME DEFAULT}, \code{SESSIONS\_PER\_USER DEFAULT} (plusieurs sessions),
        fonction forte, verrouillage court (éviter déni de service).
\end{enumerate}
\textbf{Gestion de la différence guichet (1 session) / conseiller (2 sessions) :}
deux profils distincts (\code{P\_BQ\_AGENT} et \code{P\_BQ\_CONSEIL}) car
\code{SESSIONS\_PER\_USER} est un paramètre de profil. Alternative : un seul profil +
trigger de connexion, mais moins standard et plus complexe à maintenir.
\ecart{Aucun.}

\subsection{Exercice 6.3 — T3 : Comptes}
\textbf{Objectif.} Créer GUI1, GUI2, CONS1, DIR1, AUD1, APP\_BATCH, DEV\_BQ avec le bon
profil, aucun quota, mot de passe conforme à la fonction de vérification, et le proxy
DEV\_BQ $\to$ BANQUE.

\shots{6_3d}

\textbf{Observations et réponses.}
\begin{enumerate}
  \item \code{GUI1}, \code{GUI2} (\code{P\_BQ\_AGENT}) ; \code{CONS1} (\code{P\_BQ\_CONSEIL}) ;
        \code{DIR1}, \code{AUD1} (\code{P\_BQ\_SENSIBLE}) ; \code{APP\_BATCH}
        (\code{P\_BQ\_SERVICE}) ; \code{DEV\_BQ} (\code{DEFAULT}).
  \item Aucun quota n'est affecté (la requête \code{DBA\_TS\_QUOTAS} est vide pour ces
        comptes).
  \item Les mots de passe respectent les fonctions de vérification (ex. \code{ORA12C\_VERIFY\_FUNCTION}
        rejette les mots de passe contenant le nom d'utilisateur ; plusieurs tentatives ont
        été nécessaires pour trouver des mots de passe conformes).
  \item \code{ALTER USER banque GRANT CONNECT THROUGH dev\_bq} autorise le proxy.
  \item \code{DEV\_BQ} n'a pas besoin de quota : les objets qu'il crée (index, vues)
        appartiennent au schéma \code{BANQUE} (via la connexion proxy
        \code{dev\_bq[banque]}), et \code{BANQUE} a le quota \code{UNLIMITED}.
\end{enumerate}
\ecart{Aucun.}

\subsection{Exercice 6.4 — T4 : Rôles et habilitations}
\textbf{Objectif.} Implémenter la matrice T1, attribuer les rôles aux comptes, déléguer
\code{R\_BQ\_GUICHET} à DIR1 avec ADMIN OPTION, et vérifier E1/E2.

\shots{6_4}
\shots{6_4c}

\textbf{Observations et réponses.}
\begin{enumerate}
  \item Tous les rôles sont créés avec la bonne hiérarchie d'héritage.
  \item Les rôles sont affectés aux comptes : GUI1/GUI2 $\to$ \code{R\_BQ\_GUICHET},
        CONS1 $\to$ \code{R\_BQ\_CONSEILLER}, DIR1 $\to$ \code{R\_BQ\_DIRECTEUR},
        AUD1 $\to$ \code{R\_BQ\_AUDIT}, APP\_BATCH $\to$ \code{R\_BQ\_BATCH},
        DEV\_BQ $\to$ \code{R\_BQ\_CONNEXION}.
  \item \code{GRANT r\_bq\_guichet TO dir1 WITH ADMIN OPTION} : DIR1 peut déléguer
        \code{R\_BQ\_GUICHET} (test 10 réussit) mais pas \code{R\_BQ\_CONSEILLER}
        (test 11 : \code{ORA-01932}).
  \item Contrôle E1 : tous les privilèges objet passent par les rôles (requête 6\_4c),
        aucun privilège direct aux utilisateurs.
  \item Contrôle E2 : aucun privilège système direct (\code{DBA\_SYS\_PRIVS} vide pour les
        utilisateurs métier), aucun rôle DBA/CONNECT/RESOURCE directement affecté.
  \item Arbre des rôles de DIR1 :
        \code{R\_BQ\_DIRECTEUR} $\to$ \code{R\_BQ\_CONSEILLER} $\to$ \code{R\_BQ\_GUICHET}
        $\to$ \code{R\_BQ\_CONNEXION}.
\end{enumerate}
\ecart{Aucun.}

\subsection{Exercice 6.5 — T5 : Audit}
\textbf{Objectif.} Mettre en place l'audit répondant à E6 à E8 :
\code{POL\_BQ\_CLIENTS} (modifications clients), \code{POL\_BQ\_SOLDES} (échecs seulement
sur comptes/opérations), FGA \code{FGA\_BQ\_GROS\_MONTANTS} (insertions $\ge$ 50 000),
et \code{ORA\_ACCOUNT\_MGMT} (E8).

\shots{6_5}

\textbf{Observations et réponses.}
\begin{enumerate}
  \item \code{POL\_BQ\_CLIENTS} (\code{INSERT, UPDATE, DELETE ON banque.clients}) :
        succès et échecs tracés (E6).
  \item \code{POL\_BQ\_SOLDES} et \code{POL\_BQ\_SOLDES\_OPS} activés avec
        \code{WHENEVER NOT SUCCESSFUL} : seules les tentatives de modification directe
        des soldes échouées sont tracées (E6).
  \item \code{FGA\_BQ\_GROS\_MONTANTS} sur \code{banque.operations} (\code{INSERT},
        \code{MONTANT >= 50000}) : trace les gros montants avec l'agent (E7).
  \item \code{ORA\_ACCOUNT\_MGMT} activé : trace CREATE USER, GRANT, REVOKE, DROP USER
        (E8).
\end{enumerate}
\textbf{Pourquoi POL\_BQ\_SOLDES ne trace pas les succès :} les modifications légitimes
des soldes passent par \code{p\_operation} (droits du définisseur) et sont déjà
journalisées dans la table \code{operations}. Tracer les succès noierait la trace sous du
bruit ; seules les tentatives directes et non autorisées (échecs) méritent d'être
signalées.
\ecart{\code{AUDIT\_UNIFIED\_ENABLED\_POLICIES} n'affiche pas les politiques activées
(écart récurrent 26ai).}

\subsection{Exercice 6.6 — T6 : Campagne de tests (16 tests)}
\textbf{Objectif.} Exécuter les 16 tests, comparer attendu vs obtenu, et corriger les écarts.

\shots{6_6_a}
\shots{6_6_b}
\shots{6_6_fond}

\begin{center}
\small
\begin{tabular}{|c|l|l|l|l|}
\hline
\textbf{N°} & \textbf{Compte} & \textbf{Action} & \textbf{Attendu} & \textbf{Obtenu} \\
\hline
1 & GUI1 & lire v\_clients\_guichet & succès & succès (4 lignes) \\
\hline
2 & GUI1 & lire revenu\_mensuel dans clients & ORA-00942 & ORA-00942 \\
\hline
3 & GUI1 & dépôt 2000 sur 1003 via p\_operation & succès & succès (solde 6200) \\
\hline
4 & GUI1 & UPDATE comptes solde = 1e6 & ORA-01031 + trace & ORA-41900 \\
\hline
5 & GUI1 & retrait 60000 sur 1002 via p\_operation & succès + FGA & succès + trace FGA \\
\hline
6 & GUI1 & retrait 5000 sur 1005 & ORA-20002 & ORA-20002 \\
\hline
7 & GUI1 & ouvrir 2e session & ORA-02391 & ORA-02391 \\
\hline
8 & CONS1 & modifier téléphone client 2 & succès + trace & succès + trace POL\_BQ\_CLIENTS \\
\hline
9 & CONS1 & modifier CIN client 2 & ORA-01031 & ORA-41900 \\
\hline
10 & DIR1 & accorder R\_BQ\_GUICHET à GUI3 & succès + trace & succès + trace ORA\_ACCOUNT\_MGMT \\
\hline
11 & DIR1 & accorder R\_BQ\_CONSEILLER à GUI3 & ORA-01932 & ORA-01932 \\
\hline
12 & AUD1 & lire operations + trace & succès & succès (3 opérations) \\
\hline
13 & AUD1 & DELETE FROM banque.operations & ORA-01031 & ORA-41900 \\
\hline
14 & GUI2 & 3 mauvais mdp puis le bon & ORA-28000 & ORA-28000 \\
\hline
15 & DEV\_BQ & connexion dev\_bq[banque] + index & succès & succès \\
\hline
16 & APP\_BATCH & dépôt intérêts 1200 sur 1002 & succès & succès (solde 92400) \\
\hline
\end{tabular}
\end{center}

\textbf{Écarts et corrections :}
\begin{enumerate}
  \item Tests 4, 9, 13 : l'énoncé attend \code{ORA-01031} mais 26ai renvoie
        \code{ORA-41900: missing UPDATE/DELETE privilege} (écart connu, chapitre 2).
        Le comportement est équivalent (refus d'accès) ; aucun correctif nécessaire.
  \item Test 15 : initialement échouait avec \code{ORA-01045} car \code{BANQUE} n'avait
        pas \code{CREATE SESSION}. Corrigé par \code{GRANT CREATE SESSION TO banque}.
  \item Le test 7 (2e session GUI1) a nécessité une session de fond ouverte en parallèle
        (\code{6\_6\_fond}) pour saturer la limite \code{SESSIONS\_PER\_USER 1}.
\end{enumerate}

\subsection{Exercice 6.7 — T7 : Revue des habilitations}
\textbf{Objectif.} Produire trois requêtes de revue pour l'auditeur.

\shots{6_7}
\shots{6_7b}

\textbf{Observations et réponses.}
\begin{enumerate}
  \item Requête 1 (\code{LISTAGG} des rôles par utilisateur) : 8 utilisateurs métier,
        chacun avec son profil, son état (\code{OPEN}), sa date d'expiration et ses rôles
        directs concaténés.
  \item Requête 2 : tous les privilèges objet sur le schéma \code{BANQUE} sont accordés
        aux rôles (jamais directement aux utilisateurs), conformément à E1.
  \item Requête 3 : événements d'audit par politique et utilisateur sur 7 jours
        (\code{ORA\_ACCOUNT\_MGMT}, \code{ORA\_LOGON\_FAILURES}, \code{POL\_BQ\_CLIENTS},
        \code{POL\_BQ\_SOLDES\_OPS}, \code{ORA\$MANDATORY}).
\end{enumerate}
\ecart{La requête 2 utilise \code{TABLE\_NAME} (et non \code{OBJECT\_NAME}) car
\code{DBA\_TAB\_PRIVS} n'a pas de colonne \code{OBJECT\_NAME} en 26ai.}

\subsection{Exercice 6.8 — Questions de synthèse}
\textbf{Objectif.} Répondre aux 6 questions de synthèse en prose.

\subsubsection{1. Départ de GUI2}
GUI2 quitte la banque. Procédure complète :
\begin{enumerate}
  \item \textbf{Verrouiller immédiatement} le compte :
        \code{ALTER USER gui2 ACCOUNT LOCK;} (empêche toute nouvelle connexion).
  \item \textbf{Tuer ses sessions ouvertes} :
        \code{SELECT sid, serial\# FROM v\$session WHERE username='GUI2';} puis
        \code{ALTER SYSTEM KILL SESSION 'sid,serial\#' IMMEDIATE;}
        (comme au chapitre 3.3, la révocation des rôles ne suffit pas pour les sessions
        déjà ouvertes).
  \item \textbf{Révoquer ses rôles} : \code{REVOKE r\_bq\_guichet FROM gui2;}
  \item \textbf{Vérifier qu'il ne possède aucun objet} :
        \code{SELECT owner, object\_name FROM dba\_objects WHERE owner='GUI2';}
  \item \textbf{Consulter son historique d'audit} dans \code{UNIFIED\_AUDIT\_TRAIL}
        pour garder une trace de ses dernières actions.
  \item \textbf{Supprimer le compte plus tard} (après la durée de conservation légale)
        : \code{DROP USER gui2 CASCADE;}
\end{enumerate}
\textbf{Faut-il supprimer immédiatement ?} Non. On verrouille et on révoque
immédiatement pour couper l'accès, mais on conserve le compte (et ses objets/audit)
pendant la durée de conservation requise avant suppression définitive.

\subsubsection{2. Pourquoi P\_OPERATION protège mieux qu'un GRANT UPDATE(solde)}
\code{P\_OPERATION} (droits du définisseur) centralise : contrôle de validité du montant
(\code{> 0}), vérification du solde suffisant pour les retraits (\code{ORA-20002}),
mise à jour atomique du solde \textbf{et} journalisation dans \code{operations} (agent,
montant, type, date), le tout dans une transaction avec \code{COMMIT}. Un
\code{GRANT UPDATE(solde)} direct aux guichetiers permettrait de modifier le solde
arbitrairement sans contrôle ni journalisation, contournant E3.

\subsubsection{3. Risque de la délégation du directeur}
DIR1 a \code{R\_BQ\_GUICHET} avec \code{ADMIN OPTION} : il peut créer des habilitations
excessives ou frauduleuses (accorder le rôle à n'importe qui). L'audit
\code{ORA\_ACCOUNT\_MGMT} trace chaque \code{GRANT}/\code{REVOKE} (qui, quand, quel rôle).
L'\code{ADMIN OPTION} ne se révoque pas en cascade : si on retire le rôle à DIR1, les
comptes qu'il a habilités conservent leur rôle. Il faut donc auditer et réviser
régulièrement.

\subsubsection{4. Le DBA (SYS) peut lire les revenus — acceptable ?}
Non, ce n'est pas acceptable dans un contexte strict car le DBA a un accès total aux
données sensibles (\code{revenu\_mensuel}). Solutions Oracle au-delà du TP :
Oracle Database Vault (séparation des tâches DBA / sécurité), Transparent Data
Encryption (chiffrement des données), Data Redaction (masquage dynamique), Virtual
Private Database (filtrage ligne par ligne), Oracle Label Security (étiquettes de
confidentialité), Audit Vault and Database Firewall (audit centralisé).

\subsubsection{5. Limiter le conseiller à ses clients}
Créer une vue filtrée :
\code{CREATE VIEW banque.v\_clients\_conseiller AS SELECT * FROM banque.clients WHERE
conseiller = SYS\_CONTEXT('USERENV','SESSION\_USER');} (les valeurs de la colonne
\code{conseiller} sont des noms d'utilisateur). Ou mettre en place une politique VPD
(\code{DBMS\_RLS}) qui ajoute automatiquement le prédicate
\code{conseiller = SYS\_CONTEXT('USERENV','SESSION\_USER')} sur les requêtes de
\code{clients}.

\subsubsection{6. Deux améliorations de la politique de l'agence}
\begin{enumerate}
  \item \textbf{Authentification forte / centralisée} : remplacer l'authentification
        mot de passe par Oracle Database Vault ou une authentification multifactorielle,
        et centraliser la gestion des comptes (LDAP / Active Directory) pour éviter la
        prolifération de mots de passe locaux.
  \item \textbf{Chiffrement et révision périodique} : activer TDE sur les tables
        sensibles (clients, opérations), chiffrer les échanges réseau (TLS), et mettre en
        place une revue trimestrielle des habilitations avec l'analyse des privilèges
        (\code{DBMS\_PRIVILEGE\_CAPTURE}) pour retirer les droits inutilisés.
\end{enumerate}
\ecart{Aucun.}
